\section{Criteri di valutazione}
\label{sec:background}

\subsection{Mean Dice}

Per il record di validazione $i$, $\Omega_i=\{1,\ldots,H_i\}\times
\{1,\ldots,W_i\}$ è l'insieme delle coordinate dei pixel di un'immagine alta $H_i$ e
larga $W_i$, cioè il suo dominio discreto. Siano $P_i\subseteq\Omega_i$ la maschera di
foreground predetta e $G_i\subseteq\Omega_i$ l'unione delle istanze annotate. Il
coefficiente Dice è definito come
\begin{equation}
D_i=\frac{2|P_i\cap G_i|+\varepsilon}
{|P_i|+|G_i|+\varepsilon},
\end{equation}
dove $\varepsilon$ evita una divisione per zero. Quando entrambe le maschere sono vuote,
l'implementazione assegna $D_i=1$. Sugli $N$ record annotatore-immagine del validation
set, la metrica riportata è
\begin{equation}
\mathrm{Mean\ Dice}=\frac{1}{N}\sum_{i=1}^{N}D_i.
\end{equation}
Questa misura sulla maschera unita è adatta al forte sbilanciamento tra le classi, perché
non premia i pixel di sfondo classificati correttamente. Tuttavia, non descrive come il
foreground sia diviso in oggetti distinti: un filamento frammentato in più componenti
può comunque conservare una sovrapposizione elevata a livello di pixel.

\subsection{Panoptic Quality}

La PQ viene quindi impiegata come metrica complementare a livello di istanza. Per ogni
record annotatore-immagine, le istanze predette e annotate vengono associate quando la
loro Intersection over Union è maggiore di 0.5. Con questa soglia, le associazioni valide
sono uno a uno. Le corrispondenze corrette e le istanze predette o annotate non associate
vengono accumulate sull'intero validation set come TP, FP e FN. La metrica è
\begin{equation}
\mathrm{PQ}=
\underbrace{\frac{\sum_{(p,g)\in\mathrm{TP}}\mathrm{IoU}(p,g)}{|\mathrm{TP}|}}_{\mathrm{SQ}}
\times
\underbrace{\frac{|\mathrm{TP}|}{|\mathrm{TP}|+\tfrac12|\mathrm{FP}|+\tfrac12|\mathrm{FN}|}}_{\mathrm{RQ}}.
\end{equation}
La Segmentation Quality (SQ) misura la sovrapposizione media delle istanze associate,
mentre la Recognition Quality (RQ) penalizza le istanze mancanti e quelle spurie. La
distinzione è importante in questo studio, perché configurazioni con Mean Dice simile
possono produrre numeri diversi di componenti frammentate o false positive.

Le due metriche vengono riportate per un operating point definito in modo esplicito. In
particolare, il miglior Dice di una configurazione di output non viene mai combinato con
la miglior PQ di un'altra.
